\documentclass[letterpaper]{article} %único necesario para: crear documento, hacer título, espacios, interlineado

\usepackage{vmargin} %necesario para ajustar las márgenes con \setmargins
%\usepackage{euler} %usa la fuente "euler" para las ecuaciones
\usepackage{amsfonts} %para Z estilizada
\usepackage{pifont} %para más estilos de viñetas
\usepackage{amsmath} %necesario equation*
\usepackage{amsthm} %necesario para proof enviroment
\usepackage{dsfont} %alguna letra caligráfica
\usepackage{tipa} %Epsilon bonito
\usepackage{graphicx} %para manejar imágenes
\usepackage{wrapfig} %para imágenes en modo "wrap"
\usepackage{lipsum} %para usar el texto de relleno
\usepackage{bold-extra} %permite usar textsc
\usepackage{tikz-cd} %para los diagramas conmutativos
\usepackage{mathrsfs} %para estilo especial de letra

\counterwithin{equation}{subsection}
\newtheorem{Tma}{Teorema}
\newtheorem{Prop}{Proposición}
\newtheorem{Def}{Definición}
\newtheorem{Lema}{Lema}
\newtheorem{Cor}{Corolario}
\newtheorem{Ejc}{Ejercicio}
\newtheorem{Ejm}{Ejemplo}
\newtheorem{Not}{Notación}

\newcommand{\Sh}{\normalfont{\text{Sh}}}
\newcommand{\PreSh}{\normalfont{\text{PreSh}}}
\newcommand{\PSu}{\normalfont{\text{P-Su}}}
\newcommand{\QSu}{\normalfont{\text{Q-Su}}}
\newcommand{\RSu}{\normalfont{\text{R-Su}}}
\newcommand{\Pgerm}{\normalfont{\text{P-germ}}}
\newcommand{\Qgerm}{\normalfont{\text{Q-germ}}}
\newcommand{\Rgerm}{\normalfont{\text{R-germ}}}
\newcommand{\germ}{\normalfont{\text{germ}}}
\newcommand{\Su}{\normalfont{\text{Su}}}
\newcommand{\Set}{\normalfont{\textbf{Set}}}
\newcommand{\Top}{\normalfont{\textbf{Top}}}
\newcommand{\Bund}{\normalfont{\textbf{Bund}}}
\newcommand{\Etale}{\normalfont{\textbf{Etale}}}
\newcommand{\Open}{\mathcal{O}}
\newcommand{\subseteqab}{\stackrel{ab}{\subseteq}}

\renewcommand*{\proofname}{\textbf{Prueba}} %Muestra "prueba" en lugar de "proof"
\renewcommand\refname{Referencias} %Muestra "Referencias" en lugar de "References"

%letterpaper: 21.59cm x 27.94cm
\setmargins{2.8 cm} %margen izquierdo
{1.5 cm} %margen Superior
{15.5cm} %anchura del texto
{23cm} %altura del texto
{0pt} %altura de los encabezados
{1.5cm} %espacio entre el texto y los encabezados
{5 cm} %altura del pie de página
{1 cm} %espacio entre el texto y el pie de página

\setlength{\parindent}{0pt} %Quita la sangría al iniciar los párrafos
\addtolength{\skip\footins}{2pc}

%\pagenumbering{gobble} %Quita la numeración de las páginas
\renewcommand{\baselinestretch}{1.5} %Aumenta el interlineado a n veces el automático
\relpenalty=9999 %Evita que se rompan las ecuaciones en el cambio de renglón
\binoppenalty=9999
\renewenvironment{abstract}
{\small
   \list{}{%
      \setlength{\leftmargin}{1.4cm}% <---------- CHANGE HERE
      \setlength{\rightmargin}{\leftmargin}%
   }%
   \item\relax}
{\endlist}

\renewcommand{\abstractname}{\textsc{Resumen}} %Muestra "Resumen" en lugar de "Abstract"

\usepackage{hyperref}
\hypersetup{
    colorlinks=true,
    linkcolor=blue,
    allcolors=blue,
    bookmarks=true,
    bookmarksopen=true,
    bookmarksnumbered=true
}
\usepackage{bookmark}

\date{}

\begin{document}

   \input{title.tex}
   \vspace{1cm}

   \begin{abstract}
      \textsc{Resumen.} Damos tres definiciones funtoriales de haz y probamos que son equivalentes. Luego mostramos que el prehaz de funciones reales acotadas no es un haz. 

     %   
      %\vspace{0.3cm}
      %\textsc{Palabras Clave.} Haz; espacio étalé; prehaz; homeomorfismo local; manojo; hacificación; equivalencia de categorías; local vs global.
   \end{abstract}
   \vspace{1cm}

   La Definición \ref{def:Godement} es debida a \cite[p.~109]{TATF}. 

   \begin{Def}\label{def:Godement}
      Sean $X$ un espacio topológico y $\mathscr{F}$ un prehaz de conjuntos sobre $X$. Se dice que $\mathscr{F}$ es un ${\text{haz}_1}$ de conjuntos cuando se cumplen las siguientes condiciones\footnote{Dados abiertos $U$ y $V\subseteq U$, y un $s\in\mathscr{F}(U)$, la imagen de $s$ por la aplicación $\mathscr{F}(U)\to\mathscr{F}(V)$ es la restricción de $s$ a $V$.}:
      \begin{itemize}
         \item[(F1):] Sea $(U_i)_{i\in I}$ una familia de conjuntos abiertos de $X$, $U$ la unión de los $U_i$ y $s',s''$ dos elementos de $\mathscr{F}(U)$; si las restricciones de $s'$ y $s''$ a cada $U_i$ son iguales, entonces $s'=s''$.
         \item[(F2):] Sea $(U_i)_{i\in I}$ una familia de conjuntos abiertos de $X$, con unión $U$, y supóngase dados unos $s_i\in \mathscr{F}(U_i)$ de tal manera que, cualesquiera que sean $i,j\in I$, las restricciones de $s_i$ y $s_j$ a $U_i\cap U_j$ sean iguales; entonces existe un $s\in \mathscr{F}(U)$ cuya restricción a $U_i$ es $s_i$ para todo $i\in I$.
      \end{itemize}
   \end{Def}

   \cite[p,~14]{ST} llama ``monoprehaces" o ``haces separados" a los prehaces que satisfacen \textit{(F1)}, y llama a \textit{(F2)} la ``condición de pegado". Nótese que por \textit{(F1)} se tiene que el elemento $s\in \mathscr{F}(U)$ cuya existencia garantiza \textit{(F2)} es único. Una definición alternativa de haz elimina la condición \textit{(F1)} pero exige la unicidad del $s\in \mathscr{F}(U)$ en \textit{(F2)}:

   \begin{Def}
      Sean $X$ un espacio topológico y $\mathscr{F}$ un prehaz de conjuntos sobre $X$. Se dice que $\mathscr{F}$ es una $\text{haz}_2$ de conjuntos si para cualesquiera $U\subseteqab X$, $(U_i)_{i\in I}$ cubrimiento abierto de $U$ y para cualquier colección $(s_i)_{i\in I}$ con $s_i\in \mathscr{F}(U_i)$ para cada $i\in I$, tal que\footnote{Denotamos $f^i_{ij}$ a la aplicación $\mathscr{F}(U_i)\to\mathscr{F}(U_i\cap U_j)$ y $f^j_{ij}$ a la aplicación $\mathscr{F}(U_j)\to\mathscr{F}(U_i\cap U_j)$.} $f^i_{ij}(s_i)=f^j_{ij}(s_j)$ para cualesquiera $i,j\in I$, existe un único $s\in \mathscr{F}(U)$ tal que\footnote{Denotamos $f_i$ a la aplicación $\mathscr{F}(U)\to\mathscr{F}(U_i)$.} $f_i(s)=s_i$ para todo $i\in I$. 
   \end{Def}

   La Definición \ref{def:Moerdijk} es debida a \cite[p.~66]{SGL}; para ella introducimos antes tres funciones canónicas.

   \begin{Def} 
      Sea $P$ un prehaz sobre un espacio topológico $X$. Dados $U\in\mathcal{O}(X)$ y $\left\lbrace U_i\right\rbrace_{i\in I}$ un cubrimiento abierto de $U$, definimos las siguientes funciones\footnote{Dados abiertos $U$ y $V\subseteq U$, y un $f\in PU$, representamos por $f|^U_V$ a la imagen de $f$ por la aplicación $PU\to PV$.}:
      \begin{itemize}
         \item $e:PU\to \prod_{i\in I}PU_i$ que a cada $f\in PU$ le asigna la familia $\left\lbrace f|^{U}_{U_i}\right\rbrace_{i\in I}$.
         \item $\pi_1:\prod_{i\in I}PU_i\to \prod_{(i,j)\in I\times I}P(U_i\cap U_j)$ que a cada $\left\lbrace f_i\right\rbrace_{i\in I}\in\prod_{i\in I}PU_i$ le asigna la familia $\left\lbrace {f_i}|^{U_i}_{U_i\cap U_j}\right\rbrace_{(i,j)\in I\times I}$.
         \item $\pi_2:\prod_{i\in I}PU_i\to \prod_{(i,j)\in I\times I}P(U_i\cap U_j)$ que a cada $\left\lbrace f_i\right\rbrace_{i\in I}\in\prod_{i\in I}PU_i$ le asigna la familia $\left\lbrace {f_j}|^{U_j}_{U_i\cap U_j}\right\rbrace_{(i,j)\in I\times I}$.
      \end{itemize}
      Llamamos a $e$ la función canónica de $PU$ en $\prod_{i\in I}PU_i$ y a $\pi_1$ y $\pi_2$ las funciones canónicas de $\prod_{i\in I}PU_i$ en $\prod_{(i,j)\in I\times I}P(U_i\cap U_j)$.
   \end{Def}

   \begin{Def}\label{def:Moerdijk}
      Un ${haz}_3$ de conjuntos $P$ sobre un espacio topológico $X$ es un prehaz de conjuntos tal que para cualquier $U\in \mathcal{O}(X)$ y cualquier cubrimiento abierto $\left\lbrace U_i\right\rbrace_{i\in I}$ de $U$, el siguiente es un diagrama igualador:
      \input{Diag1.tex}
   donde $e,\pi_1$ y $\pi_2$ son las respectivas funciones canónicas.
   \end{Def}
   
   Ahora queremos probar que las definiciones ${\text{haz}}_1$, ${\text{haz}}_2$ y ${\text{haz}}_3$ son equivalentes. Para ello damos primero un par de lemas. El primero de ellos nos dice que en la anterior definición todo prehaz cumple $\pi _1\circ e = \pi _2 \circ e$:

   \begin{Lema}\label{Lema:Prehaz}
      Sea $P$ un prehaz de conjuntos sobre $X$. Para cualquier $U\in \mathcal{O}(X)$ y cualquier cubrimiento abierto $(U_i)_{i\in I}$ de $U$ se tiene $\pi_1\circ e=\pi_2\circ e$, donde $e$, $\pi_1$ y $\pi_2$ son las respectivas funciones canónicas.
   \end{Lema}

   \begin{proof}
      Sea $s\in PU$. Tenemos
      $$
      \begin{aligned}
         (\pi_1\circ e)(s) &= \pi_1(e(s)) \\
                           &= \pi_1 \left(s|^U_{U_i}\right)_{i\in I} \\
                           &= \left(\left(s|^U_{U_i}\right)|^{U_i}_{U_i\cap U_j}\right)_{\left(i,j\right)\in I\times I}
      \end{aligned}
      $$
      y
      $$
      \begin{aligned}
         (\pi_2\circ e)(s) &= \pi_2(e(s)) \\
                           &= \pi_2 \left(s|^U_{U_i}\right)_{i\in I} \\
                           &= \left(\left(s|^U_{U_j}\right)|^{U_j}_{U_i\cap U_j}\right)_{\left(i,j\right)\in I\times I}
      \end{aligned}
      $$
   Para cualesquiera $i,j\in I$ tenemos en $\mathcal{O}(X)$ el diagrama\\ 
   % https://tikzcd.yichuanshen.de/#N4Igdg9gJgpgziAXAbVABwnAlgFyxMJZABgBoBGAXVJADcBDAGwFcYkQBVEAX1PU1z5CKcqWLU6TVuw4B9LDz4gM2PASIAmChIYs2iTvIA6RgMb00AAjkArRf1VCiojTqn7Dd7hJhQA5vBEoABmAE4QALZIoiA4EEjEvCHhUYgAzDRxCUkgYZFIWrHx6Tl5qYVZiOSlKQWZxYmU3EA
\begin{center}
\begin{tikzcd}
  & U_i \arrow[ld] &                                              \\
U &                & U_i\cap U_j \arrow[ld] \arrow[lu] \arrow[ll] \\
  & U_j \arrow[lu] &                                             
\end{tikzcd}
\end{center}

   de modo que en $\Set$ tenemos el diagrama\\
   % https://tikzcd.yichuanshen.de/#N4Igdg9gJgpgziAXAbVABwnAlgFyxMJZABgBoBGAXVJADcBDAGwFcYkQAFAVRAF9T0mXPkIpypYtTpNW7bgH0sfASAzY8BIgCYKUhizaJOACi6KAOuYDG9NAAIzAKwCUywepFFxWvTMOcnPikYKABzeCJQADMAJwgAWyQyEBwIJHJ+aLjExHEUtMQtTJBYhKSaVKQAZmLSnKqKgqKVOvL8pCLKXiA
\begin{center}
\begin{tikzcd}
                                    & PU_i \arrow[rd] &                \\
PU \arrow[ru] \arrow[rd] \arrow[rr] &                 & P(U_i\cap U_j) \\
                                    & PU_j \arrow[ru] &               
\end{tikzcd}
\end{center}

   Como $s\in PU$, siguiendo el diagrama tenemos
   $$
   \left(s|^U_{U_i}\right)|^{U_i}_{U_i \cap U_j}=s|^U_{U_i\cap U_j}=\left(s|^U_{U_j}\right)|^{U_j}_{U_i\cap U_j}.
   $$
   Por lo tanto $(\pi_1\circ e)(s)=(\pi_2\circ e)(s)$. Obtenemos $\pi_1\circ e=\pi_2\circ e$.
   \end{proof}

   \begin{Lema}\label{Lema:Igualadores2}
      Supongamos que en \normalfont{\textbf{Set}} el siguiente es un diagrama igualador:
      \input{Diag4.tex}
      Entonces, para todo $a\in A$ tal que $f(a)=g(a)$, existe $\alpha\in E$ tal que $e(\alpha)=a$.
   \end{Lema}
   \begin{proof}
      Definimos $F:=\left\lbrace x\in A\mid f(x)=g(x)\right\rbrace (\subseteq A)$. Se tiene que $(F,\normalfont{in}_{F,A})$ (donde $\normalfont{in}_{F,A}$ es la función inclusión de $F$ en $A$), es un igualador de $f$ y $g$, con lo cual, existe una única flecha $v:F\to E$ tal que $e\circ v=\normalfont{in}_{F,A}$. Como $a\in F$, tenemos $\alpha:=v(a)\in E$ y $e(\alpha)=e(v(a))=\normalfont{in}_{F,A}(a)=a$. 

   \end{proof}

   \begin{Prop}
      Las definiciones ${\text{haz}}_1$, ${\text{haz}}_2$ y ${\text{haz}}_3$ son equivalentes. 
   \end{Prop}

   \begin{proof}\footnote{A lo largo de la prueba, si $U$ y $V\subseteq U$ son subconjuntos abiertos, $\mathscr{F}$ es un prehaz y $s\in \mathscr{F}U$, entonces denotamos por $\left.s\right|^U_{V}$ a la imagen de $s$ por la aplicación $\mathscr{F}U\to\mathscr{F}V$.}
      \begin{itemize}
         \item[$\bullet$](${\text{haz}}_1\Rightarrow {\text{haz}}_2$) Todo ${\text{haz}}_1$ es un ${\text{haz}}_2$ por la observación realizada anteriormente (\textit{(F1)} garantiza la unicidad en \textit{(F2)}).
         \item[$\bullet$](${\text{haz}}_2\Rightarrow {\text{haz}}_3$) Supongamos que $\mathscr{F}$ es un ${\text{haz}}_2$. Sean $U\in \mathcal{O}(X)$ y $\left(U_i\right)_{i\in I}$ un cubrimiento abierto de $U$. Veamos que el siguiente diagrama en $\Set$ es un igualador:
            \begin{center}
\begin{tikzcd}
   \mathscr{F}U \arrow[r, "e"] & \prod_{i\in I}\mathscr{F}U_i \arrow[r, "\pi_2"', shift right] \arrow[r, "\pi_1", shift left] & {\prod_{\left(i,j\right)\in I\times I}\mathscr{F}(U_i\cap U_j)}
\end{tikzcd}
\end{center}

            Por el Lema \ref{Lema:Prehaz} tenemos que $\pi_1\circ e=\pi_2\circ e$. Supongamos que existen $A\in \Set$ y $u:A\to \prod_{i\in I}\mathscr{F}U_i$ en $\Set$ tales que $\pi_1\circ u=\pi_2\circ u$.
            % https://tikzcd.yichuanshen.de/#N4Igdg9gJgpgziAXAbVABwnAlgFyxMJZABgBpiBdUkANwEMAbAVxiRAB12BbOnACzgBjAE7AAYgF8AqiAml0mXPkIoAjOSq1GLNpzTDoAfWBZOWMAAIAkhM49+Q0ZKmGss+SAzY8BIgCYNanpmVkQOdn0jYE4GGAAzHAAKLFIAK05hLABzPhwASjNLK048LnhrW25eARFxCUSXU3ZBOjQLF1S89wVvZSIyVU1gnTCAQVlNGCgs+CJQOIMuJDIQHAgkdS0QtlZqBjoAIxgGAAVFHxUQWITukAWIJcRNtaQArZHwtCxDPxA9w+OZ16vjCmRyOD+IDgfCwNyecnmiw21BeiDew1Cn2+qkh+yOp3OfTC1wh1GhsIhiAAtKoEXckYgAMwo9ZPILaTFMXEAgnAy5g3K3e6PZmrVkrBjmTFQOjQqaQjFsGgTCRAA
\begin{center}
\begin{tikzcd}
\mathscr{F}U \arrow[r, "e"]               & \prod_{i\in I}\mathscr{F}U_i \arrow[r, "\pi_2"', shift right] \arrow[r, "\pi_1", shift left] & {\prod_{\left(i,j\right)\in I\times I}\mathscr{F}(U_i\cap U_j)} \\
A \arrow[ru, "u"'] \arrow[u, "v", dashed] &                                                                                              &                                                                
\end{tikzcd}
\end{center}

            Dado $t\in A$ tenemos $u(t)\ in \prod_{i\in I}\mathscr{F}U_i$, así que $u(t)$ es de la forma $(t_i)_{i\in I}$ con $t_i \in \mathscr{F}U_i$ para todo $i\in I$. Como
            $$
            (\pi_1\circ u)(t)=\pi_1(u(t))=\pi_1\left( \left( t_i\right)_{i\in I}\right)=\left( t_i|^{U_i}_{U_i\cap U_j}\right)_{(i,j)\in I\times I},
            $$
            $$  
               (\pi_2\circ u)(t)=\pi_2(u(t))=\pi_2\left( \left( t_i\right)_{i\in I}\right)=\left( t_j|^{U_j}_{U_i\cap U_j}\right)_{(i,j)\in I\times I}
            $$
            y $\pi_1\circ u=\pi_2\circ u$, entonces $t_i|^{U_i}_{U_i\cap U_j} = t_j|^{U_j}_{U_i\cap U_j}$ para cualesquiera $i,j\in I$. De este modo, por la definición de ${\text{haz}}_2$ existe un único $s_t\in \mathscr{F}U$ tal que $\left. s_t\right|^U_{U_i}=t_i$ para todo $i\in I$. Definimos $v:A\to \mathscr{F}(U):t\mapsto s_t$. Con lo anterior, para todo $t\in A$ se tiene
            $$
               (e\circ v)(t)=e(v(t))=e(s_t)=\left( \left. s_t\right|^U_{U_i}\right)_{i\in I}=\left( t_i\right)_{i\in I}=u(t),
            $$
            luego $e\circ v=u$. Si $w:A\to \mathscr{F}(U)$ es tal que $e\circ w=u$ entonces dado $t\in A$ se tiene
            $$
               \left( t_i\right)_{i\in I}=u(t)=(e\circ w)(t)=e(w(t))=\left( \left.w(t)\right|^U_{U_i}\right)_{i\in I},
            $$
            luego $\left.w(t)\right|^U_{U_i}=t_i$ para todo $i\in I$; por la unicidad de $s_t$ se tiene $w(t)=s_t=v(t)$. Por tanto $w=v$ , con lo cual $v$ es la única función de $A$ en $\mathscr{F}U$ tal que $e\circ v=u$. De lo anterior,
            \begin{center}
\begin{tikzcd}
   \mathscr{F}U \arrow[r, "e"] & \prod_{i\in I}\mathscr{F}U_i \arrow[r, "\pi_2"', shift right] \arrow[r, "\pi_1", shift left] & {\prod_{\left(i,j\right)\in I\times I}\mathscr{F}(U_i\cap U_j)}
\end{tikzcd}
\end{center}

            es un diagrama igualador y por tanto $\mathscr{F}$ es un ${\text{haz}}_3$.
         \item[$\bullet$](${\text{haz}}_3\Rightarrow {\text{haz}}_1$) Sea $\mathscr{F}$ un ${\text{haz}}_3$.
            \begin{itemize}
               \item[\textit(F1)] Sean $U\in \mathcal{O}(X)$, $\left( U_i\right)_{i\in I}$ un cubrimiento abierto de $U$ y $s',s''\in \mathscr{F}U$ tales que $\left.s'\right|^U_{U_i}=\left.s''\right|^U_{U_i}$ para todo $i\in I$. Tenemos que el siguiente diagrama es un igualador:
                  \begin{center}
\begin{tikzcd}
   \mathscr{F}U \arrow[r, "e"] & \prod_{i\in I}\mathscr{F}U_i \arrow[r, "\pi_2"', shift right] \arrow[r, "\pi_1", shift left] & {\prod_{\left(i,j\right)\in I\times I}\mathscr{F}(U_i\cap U_j)}
\end{tikzcd}
\end{center}

                  Además $e(s')=\left( \left. s'\right|^U_{U_i}\right)_{i\in I}=\left( \left. s''\right|^U_{U_i}\right)_{i\in I}=e(s'')$. Como en $\Set$ todo igualador es en particular una función inyectiva, se sigue que $s'=s''$.
               \item[\textit(F2)] Sean $U\in \mathcal{O}(X)$, $\left( U_i\right)_{i\in I}$ un cubrimiento abierto de $U$ y $\left( s_i\right)_{i\in I}$ una colección con $s_i\in \mathscr{F}(U_i)$ para todo $i\in I$, tal que $\left. s_i\right|^{U_i}_{U_i \cap U_j}=\left.s_j\right|^{U_j}_{U_i\cap U_j}$ para cualesquiera $i,j\in I$. Tenemos que el siguiente diagrama es un igualador:
                  \begin{center}
\begin{tikzcd}
   \mathscr{F}U \arrow[r, "e"] & \prod_{i\in I}\mathscr{F}U_i \arrow[r, "\pi_2"', shift right] \arrow[r, "\pi_1", shift left] & {\prod_{\left(i,j\right)\in I\times I}\mathscr{F}(U_i\cap U_j)}
\end{tikzcd}
\end{center}

                  Notemos que $\pi_1\left( \left( s_i\right)_{i\in I}\right)=\left( s_i|^{U_i}_{U_i\cap U_j}\right)_{(i,j)\in I\times I}=\left( s_j|^{U_j}_{U_i\cap U_j}\right)_{(i,j)\in I\times I}=\pi_2\left( \left( s_i\right)_{i\in I}\right)$ con $\left( s_i\right)_{i\in I}\in \prod_{i\in I}\mathscr{F}U_i$. Por el Lema \ref{Lema:Igualadores2}, existe $s\in \mathscr{F}U$ tal que $\left( s_i\right)_{i\in I}=e(s)=\left( \left.s\right|^U_{U_i}\right)_{i\in I}$, es decir, $\left.s\right|^U_{U_i}=s_i$ para todo $i\in I$. Por tanto $\mathscr{F}$ es un ${\text{haz}}_1$.

            \end{itemize}
      \end{itemize}
   \end{proof}

   \newpage
   \nocite{*}
   \bibliographystyle{apalike}
   \bibliography{refs.bib}
\end{document}
